% !TEX root = ../main.tex

\chapter{App Development and Testing}

\section{Evidence and Implementation Status}

This chapter begins with status because the supplied materials contain detailed future-tense engineering plans but no product source repository.
Table~\ref{tab:status} prevents design intent from being promoted into an implementation claim.

\begin{table}[htbp]
\centering
\caption{Current evidence status. ``Current'' refers only to the materials supplied for this thesis draft.}
\label{tab:status}
\begin{tabular}{P{4.5cm}P{3.2cm}P{6.3cm}}
\toprule
Artifact or claim & State & Basis and boundary \\
\midrule
Five-chapter research design and TC/RQ/C/E traceability & \implemented{} (document) & Controlling outline and workflow artifacts exist. \\
Interactive mobile workflow & \observed{} & HTML prototype and UI/UX mockups exist; not a production Android build. \\
Android client, backend, database, and job pipeline & \planned{} & Architecture and acceptance criteria exist; source revision was not supplied. \\
Revision-aware lineage and invalidation & \planned{} & Design specified; E1 artifacts absent. \\
Evidence-bounded question-answering effectiveness & \pending & E2 dataset, configurations, traces, and results absent. \\
Adaptive recommendation effectiveness & \pending & E3 event streams, baselines, and results absent. \\
Mobile latency and reliability & \pending & Build, device matrix, and system logs absent. \\
Five-participant prototype findings & \observed{}; \limited{} & Task counts appear in the final presentation; participant-level records absent. \\
\bottomrule
\end{tabular}
\end{table}

\section{E1: Revision and Representation Evaluation}

\subsection{Evaluation Question}

E1 tests the contribution promise that every claim-bearing artifact can be traced to the intended paper revision and that a revision update cannot silently leave stale derivatives marked current.

\subsection{Dataset and Fixtures}

The final evaluation should freeze a corpus containing:
\begin{itemize}
  \item papers with at least two known revisions;
  \item unchanged, locally edited, and structurally changed sections;
  \item single- and multi-column layouts;
  \item equations, tables, figures, appendices, and broken text layers;
  \item metadata-only or abstract-only degraded cases.
\end{itemize}
For each fixture, a gold manifest records logical identity, revision identity, expected sections, changed regions, and affected derived artifacts.

\subsection{Metrics}

E1 reports logical-identity accuracy, revision-identity accuracy, changed-section detection, stale-artifact invalidation precision and recall, section-boundary F1, reading-order error rate, and degraded-mode accuracy.
Performance must be stratified by layout and revision-change type.
Failure analysis should identify identity collision, missed revision, parse order, boundary, caption, and invalidation errors.

\begin{table}[htbp]
\centering
\caption{E1 result schema. Values must be inserted from frozen fixtures and logs.}
\begin{tabular}{P{4.6cm}P{2.2cm}P{2.2cm}P{3.8cm}}
\toprule
Measure & Baseline & MNEME & Evidence artifact \\
\midrule
Logical/revision identity accuracy & \pending & \pending & Fixture manifest + test log \\
Stale-artifact invalidation precision/recall & \pending & \pending & Dependency trace \\
Section boundary F1 & \pending & \pending & Gold section annotations \\
Degraded-mode accuracy & \pending & \pending & Parser status log \\
\bottomrule
\end{tabular}
\end{table}

\section{E2: Retrieval, Answer, and Citation Evaluation}

\subsection{Evaluation Questions}

E2 asks whether the planned section-aware, evidence-bounded pipeline improves the ability to find relevant scientific evidence and whether generated claims remain supported and sufficiently covered.
It also asks how the pipeline behaves when evidence is absent, misplaced, or surrounded by distractors.

\subsection{Frozen Scholarly Question-Answering Set}

The final dataset should contain questions written against frozen paper revisions and gold evidence spans.
It should cover factual questions, method and dataset questions, comparison questions, multi-section questions, cross-paper questions, follow-ups, and unanswerable questions.
Questions and evidence are split by paper to reduce leakage.
Each item records acceptable answer content, one or more evidence spans, answerability, and scope.

\subsection{Baselines and Ablations}

The evaluation compares:
\begin{enumerate}
  \item long-context generation without retrieval;
  \item fixed-size chunk retrieval;
  \item section-aware retrieval;
  \item section-aware retrieval with neighboring-context expansion;
  \item the full pipeline without consistency verification;
  \item the full pipeline without coverage verification;
  \item the full pipeline with identifier-only verification.
\end{enumerate}
These comparisons isolate retrieval unit, context expansion, and each verification layer.
They also guard against assuming that every additional stage improves quality or latency.

\subsection{Metrics}

Retrieval is measured with Recall@\(k\), mean reciprocal rank, and nDCG against gold spans.
Answer quality is evaluated with task-appropriate exact, semantic, or human judgments.
Citation correctness measures whether cited evidence supports a claim; citation completeness measures whether material claims are supported.
Abstention precision and recall are reported for unanswerable items.
Latency and model cost are reported separately.
Automatic evaluators are calibrated on a human-reviewed subset because reference-free evaluation is not itself ground truth~\cite{es2024ragas,saadfalcon2024ares}.

\begin{table}[htbp]
\centering
\caption{E2 result schema. Component metrics prevent a fluent answer from hiding retrieval or citation failure.}
\begin{tabular}{P{4.0cm}P{2.1cm}P{2.1cm}P{2.1cm}P{2.1cm}}
\toprule
Configuration & Retrieval & Answer & Citation & Latency / cost \\
\midrule
Long context & \pending & \pending & \pending & \pending \\
Fixed chunks & \pending & \pending & \pending & \pending \\
Section-aware & \pending & \pending & \pending & \pending \\
Full evidence-bounded pipeline & \pending & \pending & \pending & \pending \\
\bottomrule
\end{tabular}
\end{table}

\section{E3: Interest-Memory Evaluation}

\subsection{Evaluation Questions}

E3 tests whether adaptive interest state improves future paper ranking under a temporal split, whether short- and long-term state serve distinct roles, and whether users can correct the model consistently.
The evaluation does not treat historical clicks as perfect preference labels.

\subsection{Data and Baselines}

The final evaluation freezes event sequences with exposure records, event semantics, timestamps, and paper candidates.
Where real longitudinal data are unavailable, synthetic sequences may test deterministic update invariants but cannot establish user benefit.
Baselines include explicit interests only, long-term history only, recent events only, one-timescale EMA, and the full multi-timescale design.

\subsection{Metrics and Tests}

Ranking metrics include Recall@\(k\) and nDCG@\(k\) on future eligible positive events.
Diversity and novelty are reported to detect narrowing or repetition.
Adaptation delay measures how quickly a confirmed new interest affects ranking.
Negative-feedback sensitivity tests whether a bounded skip signal has an appropriate rather than overwhelming effect.
Correction consistency tests whether editing or deleting a visible interest produces predictable ranking and explanation changes.

\begin{table}[htbp]
\centering
\caption{E3 result schema. Metrics must be reported on a temporal split with exposure-aware labels.}
\begin{tabular}{P{4.0cm}P{2.1cm}P{2.1cm}P{2.1cm}P{2.1cm}}
\toprule
Model & Ranking & Diversity & Adaptation & Correction \\
\midrule
Explicit only & \pending & \pending & \pending & \pending \\
Recent only & \pending & \pending & \pending & \pending \\
Single-timescale baseline & \pending & \pending & \pending & \pending \\
Multi-timescale design & \pending & \pending & \pending & \pending \\
\bottomrule
\end{tabular}
\end{table}

\section{E4: Systems and Usability Evaluation}

\subsection{Systems Measurements}

The systems evaluation requires an immutable build, backend revision, model/provider configuration, device matrix, network conditions, and corpus snapshot.
It reports digest-load latency, paper-detail latency, Q\&A first-token and completion latency, graph expansion latency, event-upload success, synchronization conflicts, background-task completion, notification delivery, cache hit rate, and daily model cost.
Each distribution includes sample count and percentile statistics.
Failure injection covers offline operation, provider timeout, malformed document, empty retrieval, verification failure, and duplicate event upload.

\subsection{Usability Protocol}

The supplied task design is retained because it covers the core research loop:
\begin{enumerate}
  \item identify why a paper was recommended and open it;
  \item identify the main contribution and verify a source;
  \item ask a dataset question and inspect cited evidence;
  \item open the graph and inspect a related paper;
  \item change a digest preference or interest topic.
\end{enumerate}
The final dataset must contain one row per participant and task with completion, time, taps, error path, confusion code, and observation.
The study report must describe recruitment, consent, moderator instructions, prototype/build version, analysis procedure, and missing data.

\subsection{Available Formative Observation}

As stated in Chapter 2, the final presentation reports task-success counts of 5, 5, 5, 2, and 4 out of five.
Within that report, the graph task is the only major drop-off.
This observation supports the decision to make graph actions explicit.
It does not show that the subsequent redesign improves completion because no retest is supplied.

\section{Threats to Validity}

\subsection{Construct Validity}

Citation correctness is not scientific truth.
A span can support a claim in one paper while the broader literature disagrees.
Likewise, a click or long reading duration is not identical to research interest.
The evaluation therefore reports operational constructs and avoids treating them as complete measures of understanding or intent.

\subsection{Internal Validity}

Generator, embedding, and judge models can change over time.
Unrecorded provider updates may alter results.
All experiments must freeze model identifiers where possible and record dates, prompts, parameters, and outputs.
LLM-as-judge measurements require calibration against human annotations.

\subsection{External Validity}

A small course-project corpus may not represent the diversity of scientific PDFs, disciplines, or revision practices.
A five-participant formative prototype study cannot establish broad usability.
Results must be stratified and conclusions restricted to the evaluated documents, questions, devices, and participants.

\subsection{Reproducibility}

The current draft lacks the source revision, experimental data, and raw usability records needed for reproduction.
This is a present evidence gap, not a formatting issue.
The thesis can become empirically complete only when those artifacts are supplied, linked, and audited.
